\documentclass[12pt]{article}
\usepackage[utf8]{inputenc}
\usepackage[T1]{fontenc}
\usepackage{hyperref}
\title{Konkomba Language Resources for K-12 Curriculum Development}
\author{Compiled by Claude Sonnet 5.5 (claude-sonnet-5-5) for Andr\'as Kornai}
\date{}
\begin{document}
\maketitle

\section{Introduction}
Konkomba (Likpakpaanl, ISO 639-3 xon) is a Gur language of northern Ghana and northern Togo. Its speakers call themselves Bikpakpaam. Speaker counts in the sources range from about 920,000 to 1.29 million. Dialect names vary between sources.

\section{Primary Educational Resources}
No primers or textbooks were found in the search results.

\section{Linguistic and Reference Materials}
\begin{itemize}
\item Likpakpaanl--English dictionary on Webonary (SIL), compiled by Mary Steele of GILLBT, fourth edition, web edition 2015: \url{https://www.webonary.work/likpakpaanl/about/introduction/}
\item Wikipedia, Konkomba language: \url{https://en.wikipedia.org/wiki/Konkomba_language}
\item Glottolog reference record: \url{https://glottolog.org/resource/reference/id/sil16:35165}
\end{itemize}

\section{Religious and Community Resources}
\begin{itemize}
\item Scripture Earth page for Konkomba, with links to Bible editions: \url{https://scriptureearth.org/00eng.php?iso=xon}
\item Joshua Project language page: \url{https://joshuaproject.net/index.php/languages/xon}
\item Wikipedia, Konkomba people: \url{https://en.wikipedia.org/wiki/Konkomba_people}
\end{itemize}

\section{Digital and Audio Resources}
\begin{itemize}
\item Omniglot page on Konkomba: \url{https://www.omniglot.com/writing/konkomba.htm}
\item Ghana Institute of Linguistics, Literacy and Bible Translation (GILLBT) page on the Webonary site: \url{https://www.webonary.work/likpakpaanl/?p=23}
\end{itemize}

\section{Curriculum Development Context}
Konkomba has a bilingual dictionary and a full Bible translation, which provide a base for orthography and vocabulary. No school materials were found in the search. GILLBT is the organisation connected with the language work and is the logical contact for literacy materials.
\end{document}